\documentclass[11pt,letterpaper]{article}

\usepackage[T1]{fontenc}
\usepackage[utf8]{inputenc}
\usepackage{geometry}
\usepackage{hyperref}
\usepackage{longtable}
\usepackage{microtype}
\usepackage{parskip}

\geometry{margin=1in}

\title{Hyperledger Cacti Whitepaper\\Version 2}
\author{Hyperledger Cacti Contributors}
\date{\today}

\begin{document}

\maketitle

\begin{abstract}
Hyperledger Cacti provides a framework for interoperability across distributed
ledger technologies and related enterprise systems. This second whitepaper is a
living source document for the architecture, trust assumptions, deployment
models, and research foundations of the project.
\end{abstract}

\tableofcontents
\newpage

\section{Status}

This document is the version 2 source scaffold. It intentionally starts as a
concise outline so maintainers can evolve the whitepaper in reviewable changes.

\section{Motivation}

Enterprise interoperability requires systems that can coordinate actions,
state, identities, and evidence across networks with different governance
models. Cacti brings together connector, routing, validation, and application
components so deployments can integrate ledgers without treating every
interoperability problem as a one-off integration.

\section{Scope}

The whitepaper should describe:

\begin{itemize}
  \item Cacti's goals and non-goals.
  \item Supported interoperability patterns.
  \item Core architectural components and extension points.
  \item Security and trust assumptions.
  \item Operational and deployment models.
  \item Related research, specifications, papers, and theses.
\end{itemize}

\section{Architecture Outline}

\subsection{Connectors}

Connectors expose ledger-specific capabilities through common APIs and plugin
interfaces.

\subsection{Validation and Verification}

Validation components provide evidence handling, proof verification, and
cross-network state checking where supported by the target networks.

\subsection{Applications and Tooling}

Cacti includes application-level packages, examples, and operational tooling
that demonstrate how interoperability components are composed.

\section{Research and Prior Art}

The bibliography is intentionally kept in source form so papers, theses, and
specifications can be reviewed alongside the text that cites them.

\nocite{hyperledger-cacti-docs}
\nocite{hyperledger-cacti-repo}

\bibliographystyle{plain}
\bibliography{references}

\end{document}
